\section{An exact local change of scale}
\label{sec:blocking}

We now construct the transformation of interactions that will be iterated.
Its defining identity holds for each realization of the disorder.  The
probability estimates needed to control its iteration enter only in
Section~\ref{sec:norms}.  The distinction matters: a principal interaction
can be large, and the construction must still preserve the normalized
partition function.  We integrate such interactions through positive
conditional weights, while expanding only the remaining small interactions.
The connected expansion used below is the list construction of
\cite[Section~5, Proposition~5.2]{ReferenceIsing}, with this additional
positive-weight operation; we include its algebra and counting estimates.

\subsection{Lists and the two kinds of support}

Let $r_0$ and $L$ be odd positive integers and set
\begin{equation}
 r_j=r_0L^j,\qquad r=r_j,\qquad R=Lr.
 \label{block:scales}
\end{equation}
Partition $\Z^2$ into disjoint squares of side $r$, with centers in
$r\Z^2$, and denote this collection by $\cG_r$.  These partitions are
nested, and the origin is a cell center at every scale.  Cells are adjacent
when their closures meet, so corner adjacency is allowed.  An
\emph{animal} is a nonempty finite connected set of cells.  Write
$\operatorname{par}(c)$ for the unique $R$-cell containing the $r$-cell $c$.
Distances between cells, when expressed without units, are maximum-norm
cell distances.  Physical distances carry the factor $r$ or $R$ explicitly.

At scale $r$ the index set is
\[
 \mathcal I_r=\{p_c:c\in\cG_r\}
 \mathbin{\dot\cup}
 \{(A,c): A\subset\cG_r\text{ is an animal},\ c\in A\}.
\]
The index $p_c$ is the \emph{principal slot} at $c$.  It may use the cells
within a fixed distance $S$ of $c$, called its halo.  The slot $(A,c)$ is
an \emph{animal slot}, with anchor $c$.  Assign sizes
\[
 n(p_c)=1,\qquad n(A,c)=|A|.
\]
Thus the principal size is one even though its allowed support has a
fixed halo.  A list $h=(h_i)_{i\in\mathcal I_r}$ consists of real functions
of finitely many spins.  Its background functions are invariant under a
simultaneous reversal of all spins.  Terms with an identical slot index
are added as functions before a norm is taken.  Different list
representations of the same total interaction need not coincide.

Every slot records two supports.  Its \emph{spin carrier} contains the
spin coordinates used by its function.  Its \emph{dependence carrier}
contains the microscopic edge variables and decision seeds on which that
function, including its choice of representation, depends.  Both carriers
are contained in its animal, or in the permitted halo of a principal
slot.  In every geometric contact test and output-footprint rule below,
we use this full indexed animal or full principal halo, even if the
actual spin or random dependence is smaller.  These geometric carriers
therefore depend only on the slot indices.  An edge variable is assigned to its incident cells; enlarging a
carrier by one safety cell removes ambiguity at interfaces.  Consequently
sets of slots whose enlarged dependence carriers are at distance greater
than two use disjoint primitive random variables and are independent.
This is an assertion about the disorder probability space, not about
spins in the reference Gibbs measure.  Independent seeds are reserved at
all cells and all future scales.  A decision at a given scale reads only
the seeds and edge variables in its recorded dependence carrier.

In the plane, the law of the list is invariant under translations of its
grid and the symmetries of the square.  Geometric choices of cutting squares are equivariant random choices
using these independent seeds.  For example, a rounding tie is resolved
by assigning independent continuous priorities to its candidates and
choosing the smallest priority.  We make one realized choice before
performing conditional averaging or taking a logarithm.  Assigning a
resulting function to an anchor is a linear operation; only for these
assignments may we instead split the function equally among all tied
slots.  Both conventions preserve square symmetries in law.

Let $\langle\cdot\rangle_0$ denote critical plane reference expectation,
acting on spins only, with the disorder fixed.  We use the projection
\begin{equation}
 \pi f=f-\langle f\rangle_0,
 \qquad \|\pi f\|_\infty\le 2\|f\|_\infty.
 \label{block:centering}
\end{equation}
All list entries are plane-centered.  The removed scalar is saved
separately, with its dependence carrier, and never subsequently treated
as an interaction.  This convention will also be applied to coefficients
containing source labels.  It is essential to save their scalars: only
source-independent scalars cancel automatically in a normalized partition
function.

For a finite set of marked sites, introduce commuting variables
$\zeta_1,\ldots,\zeta_m$ satisfying $\zeta_a^2=0$.  A \emph{source jet}
is an expression
\begin{equation}
 H(\zeta)=\sum_i h_i+
       \sum_{\varnothing\ne J\subset\{1,\ldots,m\}}
             \zeta_J\sum_i u_i^J,
 \qquad \zeta_J=\prod_{a\in J}\zeta_a.
 \label{block:jet}
\end{equation}
A coefficient carrying $J$ must have a carrier containing the cells of
all the corresponding marks.  A one-mark list is anchored at its mark
cell and has a marked principal slot there.  For more than one mark we
use animal slots and anchor at the first mark, in the prescribed order
of labels.  Source coefficients are not used in threshold decisions.
The nilpotent variables record derivatives and products exactly; they
place no positivity assumption on the source coefficients.  We always
use finitely many labels and source lists with finite weighted indexed
norms, defined below.

\subsection{Indexed norms and the order of the parameters}

The function space at a slot is the finite-dimensional space of functions
of its available spins, with its sup norm.  For a tensor in $k$ such
spaces we use the projective norm
\begin{equation}
 \|F\|_\pi=
 \inf\left\{\sum_{\alpha}\prod_{b=1}^k
            \|f_{b,\alpha}\|_\infty:
       F=\sum_{\alpha}f_{1,\alpha}\otimes\cdots\otimes f_{k,\alpha}
       \right\}.
 \label{block:projective}
\end{equation}
The infimum is over finite tensor representations, or their norm
completion.  In particular,
$\|f_1\otimes\cdots\otimes f_k\|_\pi\le\prod_b\|f_b\|_\infty$.
A bounded multilinear operation on the component function spaces extends
to this tensor product with the same operator bound.  Expectations of
random tensors are taken componentwise; the inequality
$\|\E F\|_\pi\le\E\|F\|_\pi$ then applies.

For $v>0$ and $1\le p<\infty$, define the indexed norms
\begin{align}
 \|h\|_{v,\infty}
   &=\sup_{i\in\mathcal I_r}e^{vn_i}\|h_i\|_\infty,
 \label{block:det-indexed}\\
 \|h\|_{v,p}
   &=\sup_{i\in\mathcal I_r}e^{vn_i}
                      \bigl(\E\|h_i\|_\infty^p\bigr)^{1/p},
 \label{block:indexed}\\
 \|K\|_{v,\pi}
   &=\sup_{i_1,\ldots,i_k}
          e^{v(n_{i_1}+\cdots+n_{i_k})}
                       \|K_{i_1,\ldots,i_k}\|_\pi.
 \label{block:tensor-indexed}
\end{align}
The first norm is deterministic; whenever a random essential supremum
is intended it will be stated.  In the tensor norm the factors retain
separate spin variables, even if their carriers overlap.  Diagonal
evaluation and multiplication are bounded operations, but do not replace
the projective norm.  Source norms use their marked index sets.  Rooted
sums that arise when slots are combined will be bounded by the counting
argument below; they are not part of the definition of these indexed
suprema.  An admissible source jet, at a fixed disorder realization,
has finite deterministic norms $\|u^J\|_{v,\infty}$ at a fixed weight
with counting slack, for each of its finitely many nonempty label sets.  A finite-volume list
also has only finitely many possible anchors.  These conditions imply
absolute summability, as verified in Proposition~\ref{block:finite},
and hold for the microscopic initialization and every finite iterate.
On the infinite plane we require only the corresponding indexed and
rooted estimates.

Here is an order of choices sufficient for all later uses.  Fix the
finite set of even source powers, including $2$, $4$, and $400$, and
then a much larger even background power $p$.  Choose a fixed contour
dilation $H$ large enough for the buffered reference estimates and the
finite testing statement in Section~\ref{sec:reference}.  Choose the
principal halo $S\ge 10^5H$, increasing its absolute multiplier if a
larger protective neighborhood is used.  Let $a_{\rm an}$ be an
exponential counting constant for rooted animals, established below.
One may take a slack weight $\omega>100a_{\rm an}$, a tensor weight
$10\omega$, and background and moment weights $w,v$ larger than
$10^6(p+400)^4$ times that tensor weight.  Only finitely many intermediate
weights, all with this counting slack, will occur.  They are fixed before
$L$.  Set $M=(\log L)^3$ and increase the odd integer $L$ until it pays
all the weight comparisons below and satisfies
\begin{equation}
 S\ll M,\qquad 1000M\ll L.
 \label{block:M}
\end{equation}
Choose $\delta>0$ sufficiently small after $L$.  The testing accuracy
needed later can then be chosen at this fixed $L$, followed by a
sufficiently large odd $r_0$.  Finally the disorder strength is made
small after these choices.  Constants $C_L$ may depend on the fixed
weights and geometric constants, but not on the scale, mesh, volume,
or principal field amplitudes.  The counting constants $C_L$ below can
be chosen polynomial in $L$ for fixed weights.

The deterministic hypothesis on the nonprincipal background is
\begin{equation}
 \|h_{A,c}\|_\infty\le\delta e^{-w|A|}.
 \label{block:tail-hypothesis}
\end{equation}
There is no such bound on principal slots.  At a fixed finite-volume
step they are finite real functions, with a possibly large deterministic
bound.  No stochastic moment hypothesis is used in this section.

\subsection{Positive treatment of the large principal fields}

For each principal slot choose an independent threshold
$\theta_c$ uniformly in $[\delta,2\delta]$, independently of the input
list, using a fresh seed.  Declare $c$ bad when
$\|h_{p_c}\|_\infty>\theta_c$.  All other entries are called small
entries; in particular their principal norms are at most $2\delta$.
Put an edge between bad centers at physical maximum-norm distance at
most $100HR$.  A component is selected if its diameter is at most
$1000Mr$.  The empty set has no selected component.

Selection can be decided locally.  Starting at a bad center $x$, inspect
all bad centers within distance $100HR+1000Mr$ of $x$, and the component
of $x$ in this finite inspected graph.  If this component contains a
vertex at distance greater than $1000Mr$ from $x$, the true component
is already too large.  Otherwise every possible first edge leaving the
inspected component has its other endpoint in the inspected ball, so
there is no such edge.  We have found the entire component and can test
its diameter.  This also rejects every infinite component using a
finite certificate.  The badness of each inspected center uses only
its principal halo and its fresh threshold seed.  Thus selection has a
bounded dependence neighborhood of radius $C HR$.

For a selected component $\mathcal C$, choose a lattice center near the
center of its bounding box, resolving rounding symmetrically, and let
$Q_{\mathcal C}$ be the cutting square of radius $HR$ there.  Write
\[
 B_{\mathcal C}=\sum_{c\in\mathcal C}h_{p_c},\qquad
 A_{\mathcal C}=P_{Q_{\mathcal C}}e^{B_{\mathcal C}},\qquad
 b'_{\mathcal C}=\log A_{\mathcal C}.
\]
Here and throughout, $P_Q$ is conditional expectation in the critical
reference measure given the spins outside the interior of $Q$, as
defined in Section~\ref{sec:reference}.  In particular it retains the
cutting spins.  Our choices of $L$ ensure that all spin supports of
$B_{\mathcal C}$ lie strictly inside $Q_{\mathcal C}$.  Distinct selected
squares have disjoint interiors and large intervening collars: their
bad centers belong to distinct components and hence are separated by
more than $100HR$, whereas their diameters are $O(Mr)$.
An unselected bad insertion also has its entire spin support outside
these collars.  Since the integrand is positive,
\begin{equation}
 e^{-\|B_{\mathcal C}\|_\infty}
 \le A_{\mathcal C}\le e^{\|B_{\mathcal C}\|_\infty},
 \qquad
 \|\pi b'_{\mathcal C}\|_\infty
       \le2\sum_{c\in\mathcal C}\|h_{p_c}\|_\infty.
 \label{block:positive-log}
\end{equation}
Put $\pi b'_{\mathcal C}$ in a nearest parent principal slot and save
its scalar.  For every unselected bad input, retain $h_{p_c}$ without
averaging and add it to the principal slot at $\operatorname{par}(c)$.
A boundary modification, specified below, may also retain a selected
component; its individual insertions are then treated this same way.
The pure output of all bad backgrounds is therefore principal.

For any spin function $F$, the normalized tilted conditional operator is
\begin{equation}
 P_Q^B F=\frac{P_Q(e^BF)}{P_Qe^B}.
 \label{block:tilt}
\end{equation}
It is defined for every real bounded $B$, sends $1$ to $1$, and obeys
$\|P_Q^B F\|_\infty\le\|F\|_\infty$.  These properties require no
smallness or ferromagnetic property of $B$.  The operator still uses
only reference conditional probabilities.  The denominator in
\eqref{block:tilt}, and its logarithm already saved above, are the
normalization that will make the ensuing partition identity exact.

\subsection{The transformed moments and their logarithm}

Keep the bad configuration and its selected squares fixed.  Expand the
exponential of all small entries and source coefficients in ordered
occurrences, with the factor $1/k!$ for $k$ occurrences.  An occurrence
remembers its input slot and its source label set, if any.  Retain only
products whose source label sets are disjoint, as required by
$\zeta_a^2=0$.  Repeated background indices are allowed and their
occurrences remain separately labeled.

For each such list, join two occurrences when their full geometric
carriers are within $1000HR$.  Call this contact graph $G$.  First average the product of the occurrences jointly inside
every selected square, using $P_{Q_{\mathcal C}}^{B_{\mathcal C}}$.
These operators commute because their interiors are disjoint and the
reference specification has the nearest-neighbor Markov property.
An additional ordinary conditional average is used for an occurrence
if it is a singleton component of $G$, has $n_i<L/2$, and has no input
bad center within $1000HR$ of its carrier.  The additional cutting square
has radius $HR$ and center the center of the parent of its anchor.
Its support is inside that square, and its collar is disjoint from all
other occurrence supports and all bad supports.  The ordinary cuts
are also mutually separated.  Long singletons, and short singletons
near bad centers, receive only the selected-square operations already
specified.  This defines a transformed moment $m(I)$ for every finite
labeled set $I$ of occurrences; set $m(\varnothing)=1$.

For clarity, let $\mathcal J$ index the small background entries and all
source entries.  An index $\iota\in\mathcal J$ carries a function
$f_\iota$ and a label set $J_\iota$, empty for a background entry.  Write
$m(f_{\iota_1},\ldots,f_{\iota_k})$ for the transformed moment just
constructed.  Thus its arguments already include the input functions.
The transformed series is
\begin{equation}
 \mathcal F(\zeta)=\sum_{k\ge0}\frac1{k!}
    \sum_{\substack{\iota_1,\ldots,\iota_k\in\mathcal J\\
                  J_{\iota_a}\cap J_{\iota_b}=\varnothing\ (a\ne b)}}
       \zeta_{J_{\iota_1}}\cdots\zeta_{J_{\iota_k}}
                    m(f_{\iota_1},\ldots,f_{\iota_k}),
 \label{block:moment-series}
\end{equation}
where $\zeta_\varnothing=1$ and the term $k=0$ is $1$.
Its jointly degree-zero term in all occurrence amplitudes is $1$.
Define its connected coefficients by the ordinary formal logarithm.
More explicitly, for $k$ labeled occurrences,
\begin{equation}
 c(I)=\sum_{\mathcal P\text{ partition of }I}
       (-1)^{|\mathcal P|-1}(|\mathcal P|-1)!
                          \prod_{J\in\mathcal P}m(J).
 \label{block:partition-coefficient}
\end{equation}
The ordinary variables used to distinguish occurrences in this formula
can equivalently be made square-free.  The factor $1/k!$ is retained
when these labeled coefficients are summed over ordered background
occurrences.  Lemma~\ref{block:tree} below proves convergence on the full
required ball of small entries, uniformly over the fixed large fields.
It consequently defines a real logarithmic interaction there, together
with a finite jet in the source variables.

For $|I|=1$ and $n_i<L/2$, with at most one source label, assign the
centered output to the principal slot at the parent anchor, also for
a singleton mark.  The same assignment applies when the short entry has not been averaged because
of a nearby unselected bad field.  All long singleton outputs, all
outputs carrying more than one mark, and all outputs with at least
two occurrences are put in animal slots.  For these animal slots, take the union of the parent footprints of the
input geometric carriers and thicken it by $10000H$ parent cells.  This
union is connected whenever $c(I)$ can be nonzero, as proved below;
otherwise its coefficient is zero.  Increasing this fixed thickening
if necessary includes the cutting contours and every decision
neighborhood.  Average the anchor over the parents of the input
anchors, with multiplicity, and keep the stated first-mark anchor for
marked coefficients.  Center each resulting function and save its
scalar with the same carrier before discarding its spin support.

At zero background all principal inputs are below threshold.  The
first derivative on backgrounds and one-mark lists is therefore the
operator $T$ given by
\begin{equation}
 (i,f)\longmapsto
 \begin{cases}
  \bigl(p_{\operatorname{par}(c)},\pi P_{HR,\operatorname{par}(c)}f\bigr),
       &n_i<L/2,\\
  \text{the same function in its parent animal slot},&n_i\ge L/2.
 \end{cases}
 \label{block:linear}
\end{equation}
Here the first cutting square is centered at the indicated parent
center.  For a spin source whose scaling dimension is $d=1/8$, the
scaled one-mark linear transport is $L^dT$; with $q$ spin labels the
corresponding rescaling factor is $L^{qd}$.  These rescalings are a
choice of coordinates on sources, after the exact unscaled identity.

\begin{lemma}[Carriers after one step]
\label{block:carrier}
The construction has exact finite dependence carriers.  Its principal
outputs fit in a fixed halo $S$ at scale $R$.  An animal output generated
by input occurrences of sizes $n_1,\ldots,n_k$ has size
\begin{equation}
 n'\le\sum_{a=1}^k\left(C+\frac{Cn_a}{L}\right),
 \label{block:compression}
\end{equation}
where $C$ depends only on the fixed protective neighborhoods and not on
$L$.  Both the spin carrier and the dependence carrier satisfy these
statements.  The same bound holds for the footprint of a saved scalar.
\end{lemma}

\begin{proof}
A connected animal of $n$ fine cells has a spanning-tree traversal of
length at most $2n$ in fine-cell units.  Divide this walk into consecutive
pieces of length at most $L$.  Each piece, together with its prescribed
$O(R)$-thickening, lies in a bounded number of parent cells.  There are
at most $1+2n/L$ pieces, proving the bound for one carrier.  A principal
halo has fine diameter $O(S)$ and, since $L$ is large compared with $S$,
its parent thickening has bounded size.  Taking the union proves
\eqref{block:compression}.  The $1000HR$ contact rule makes the thickened
union connected.  Every cut and every tested bad center is within a
fixed $HR$ distance of one of the input carriers.  Their input principal
halos add only $Sr\ll R$.  The selection test and fresh rounding seeds
are within the same neighborhood.  This proves the dependence claim,
not just the assertion about spin coordinates.  In a pure selected bad
output the whole component lies in an $O(Mr)$ square; its cuts and tests
thus fit in a fixed parent halo.  An unselected bad input uses an even
smaller halo.  Short singleton outputs fit in a fixed parent halo for
the same reason.  Choosing $S$ to contain these fixed neighborhoods
completes the proof.
\end{proof}

\subsection{Exactness and the connected expansion}

Let $\mu_0$ be a finite-volume reference measure in which all the cuts
just described are legal.  Let $B_{\rm ret}$ be the sum of the retained
bad insertions and set
\[
 B'=B_{\rm ret}+\sum_{\mathcal C\text{ selected}}b'_{\mathcal C}.
\]
The next statement includes all scalars removed from $B'$ and from
$\log\mathcal F$.

\begin{proposition}[Exact partition identity]
\label{block:exact}
For sufficiently small $\delta$ at fixed $L$, let a finite-volume real
background have finitely many possible anchors, satisfy
\eqref{block:tail-hypothesis}, and have finite principal norms.  Every
admissible source jet with finitely many labels then has a transformed
centered jet $H'(\zeta)$ and a saved scalar jet $s'(\zeta)$ such that
\begin{equation}
 \int e^{H(\zeta)}\,\dd\mu_0
       =e^{s'(\zeta)}\int e^{H'(\zeta)}\,\dd\mu_0.
 \label{block:partition-identity}
\end{equation}
The equality is coefficientwise in the square-free variables and holds
for every disorder realization and every choice of decision seeds.
Consequently the normalized source partition function is
\begin{equation}
 \frac{\int e^{H(\zeta)}\,\dd\mu_0}{\int e^{H(0)}\,\dd\mu_0}
 =e^{s'(\zeta)-s'(0)}
   \frac{\int e^{H'(\zeta)}\,\dd\mu_0}{\int e^{H'(0)}\,\dd\mu_0}.
 \label{block:normalized-identity}
\end{equation}
All pure bad output is principal, and the linearization at zero is
\eqref{block:linear}.
\end{proposition}

\begin{proof}
Fix one product $F$ of expansion occurrences.  Conditional on all
selected cutting spins and all remaining exterior spins, the reference
law factors over the selected interiors.  The retained bad supports
are outside these interiors.  Therefore
\begin{equation}
 \int e^{B_{\rm ret}+\sum_{\mathcal C}B_{\mathcal C}}F\,\dd\mu_0
 =\int e^{B'}\left(
       \prod_{\mathcal C}P_{Q_{\mathcal C}}^{B_{\mathcal C}}
                         \right)F\,\dd\mu_0.
 \label{block:tilted-identity}
\end{equation}
This is the tower property with the factors $P_Qe^B$ left outside the
normalized operators.  If an ordinary singleton cut is used, its
interior has a collar from all other factors, including the selected
normalizations and retained bad fields.  Conditional averaging of its
one occurrence consequently leaves the right-hand integral unchanged.
The resulting integrand is precisely $e^{B'}m(I)$.  Thus each expansion
term preserves its integral.  Summing first for finitely many inputs
and sufficiently small scalar amplitudes gives
$\int e^{H(\zeta)}\dd\mu_0=\int e^{B'}\mathcal F(\zeta)\dd\mu_0$.
By construction $B'+\log\mathcal F=s'+H'$.  Lemma~\ref{block:tree}
justifies this logarithm, its exponentiation, and the extension to the
whole small-entry ball, including finite source jets.  Absolute
finite-volume summability, proved in Proposition~\ref{block:finite},
permits truncation of an infinite slot list.  This proves
\eqref{block:partition-identity}.  Its denominator at zero sources is
strictly positive, so division gives
\eqref{block:normalized-identity}.  No disorder expectation was taken.
\end{proof}

\begin{lemma}[Tree coefficients and weighted summability]
\label{block:tree}
Fix an arbitrary real bad background and its selected squares.  For
$k$ small or marked occurrences with contact graph $G$, their connected
coefficient before centering satisfies
\begin{equation}
 c(I)=0\quad\text{if }G\text{ is disconnected},\qquad
 \|c(I)\|_\infty\le
       2^{k-1}\tau(G)\prod_{a=1}^k\|f_a\|_\infty
       \quad\text{otherwise},
 \label{block:tree-coefficient}
\end{equation}
where $\tau(G)$ is the number of spanning trees and $\tau(G)=1$ for a
single vertex.  For any of the fixed output weights, sufficiently large
$L$ allows summation of the ordered order-$k$ coefficients, with their
$1/k!$, at a fixed output anchor at cost $C_L^k$ times the product of
the input weighted bounds.  This holds also with a fixed larger output
weight, before or after centering, and for a fixed finite number of
marked occurrences or substitutions of input norm factors.

In particular the bounds can be realized by deterministic nonnegative
coefficient majorants that retain every separate factor
$\|f_a\|_\infty$.  If all but one of these factors are small backgrounds,
summing terms of order at least two leaves a factor at most $C_L\delta$
times the retained factor's weighted input bound.  This last statement
also holds with a marked retained factor under the normalized bad
operators.  The analogous multilinear maps on indexed projective
tensors have the same bounds, after the indicated allocation of
weights.
\end{lemma}

\begin{proof}
For any subset of labeled occurrences, its transformed moment factors
as a function over the components of its induced contact graph.
Indeed a selected cutting square cannot meet supports in two different
components: their distance would then be at most $2HR$ plus their
halos, much smaller than the contact threshold.  Each normalized
operator therefore acts on only one component and sends empty factors
to $1$.  Ordinary singleton cuts likewise act on just that component.
This proves algebraic factorization without asserting probabilistic
independence of the spin functions.

Regard each nonempty connected subset of labels as a polymer with
activity its transformed moment.  Two polymers are incompatible if
they overlap or a contact edge joins them.  The hard-core polymer
partition polynomial equals the transformed moment polynomial:
every induced contact graph has a unique decomposition into connected
components.  Its logarithmic coefficient is the sum over partitions
of the labels into connected blocks, with the hard-core Mayer
coefficient of the quotient contact graph.  For a finite graph $Q$
that coefficient is
\[
 \sum_{F\subset E(Q):\,(V(Q),F)\text{ connected}}(-1)^{|F|}.
\]
Its absolute value is at most $\tau(Q)$.  To see this directly, strictly
order the edges and group each connected subgraph by its greedy
minimum spanning tree.  For a specified tree, an additional edge is
permitted exactly when it is later than every edge on its tree path.
Such edges are optional independently, so their alternating sum is
zero or has absolute value one.  The sum over trees gives the bound.

For every connected block choose an internal spanning tree, and for
every edge of a quotient spanning tree choose one original contact
edge between the corresponding blocks.  Allowing all these choices
can only enlarge the bound.  Their union is a spanning tree of $G$
with each edge marked internal or between blocks.  Conversely these
marks recover the blocks and all the chosen edges.  There are at most
$2^{k-1}$ markings per spanning tree.  Every polymer activity has norm
at most the product of its occurrence norms, since all its conditional
operators, tilted or ordinary, are sup-norm contractions.  This proves
\eqref{block:tree-coefficient}, including the vanishing for disconnected
$G$.  Labeling occurrences has avoided any difficulty with repetitions;
returning to the exponential series gives the stated $1/k!$.

We give the placement count because the indexed norms alone contain
no summation.  A rooted animal of $n$ cells has a depth-first traversal
of a spanning tree with at most $2n$ nearest or corner steps.  Choosing
a deterministic traversal for counting gives at most $8^{2n}$ possible
encodings.  Choices of an anchor and of a distinguished contact cell
contribute polynomial factors in $n$, hence can be included in
$e^{a_{\rm an}n}$ after enlarging $a_{\rm an}$.
If a parent occurrence has size $n_1$, a child occurrence of size $n_2$
touching its $1000HR$ neighborhood can be placed in at most
\begin{equation}
 C L^2(1+n_1)e^{a_{\rm an}n_2}
 \label{block:child-count}
\end{equation}
ways.  One chooses a fine cell in the parent's neighborhood, at cost
$CL^2(1+n_1)$, and then the child's animal, its contact cell and anchor.
Principal halos change only the constant in this estimate.  A root
whose parent anchor is a fixed cell costs at most $CL^2$ anchor choices
and the same exponential animal count.  The average assignment of
anchors permits rooting at an input whose parent is the required
output anchor; summing this choice costs at most $k$, harmless below.
For a marked output, root instead at the occurrence containing its
first mark.

By Lemma~\ref{block:carrier}, an output weight $v_{\rm out}$ costs
\[
 e^{v_{\rm out}n'}\le
       e^{Cv_{\rm out}k}
           \prod_{a=1}^ke^{Cv_{\rm out}n_a/L}.
\]
Choose $L$ so that $Cv_{\rm out}/L$ uses only a small fixed fraction
of the smallest available input weight.  The remaining weight pays
the animal count and leaves positive slack $b$ at each vertex.  At a
vertex with $q$ children, the factors $1+n$ in
\eqref{block:child-count} are paid by
\begin{equation}
 (1+n)^q e^{-bn}\le q!\,C_b^{q+1},\qquad n\ge1.
 \label{block:factorial-slack}
\end{equation}
For example, apply $x^q/q!\le e^x$ to a fixed fraction of $bn$,
and absorb the replacement of $n$ by $1+n$ in $C_b$.
Thus a rooted labeled tree contributes at most
$C_L^k\prod_a q_a!$ times the input weighted bounds.  Ordering the
children at every vertex introduces exactly $\prod_a q_a!$ orders.
After division by $k!$, these objects are counted by plane rooted
trees, whose number at size $k$ is at most $4^{k-1}$.  Factors $k$
and $2^{k-1}$ are absorbed by increasing $C_L$.  This proves the
claimed geometric order bound.  Only $L^2$ covering factors have
entered, so $C_L$ is polynomial in $L$ for fixed weights.

The argument kept the norm of every occurrence as a separate
nonnegative factor.  The assigned output footprint depends only on
the input indices, and the anchor assignment has at most $k$ choices,
with only a fixed additional multiplicity for parent rounding.
We dominate these allowed assignments by their sum.  Conditional
kernels are bounded uniformly by one; we do not sum over badness
patterns or possible cutting kernels.  This gives deterministic
nonnegative kernels independent of the bad amplitudes.  One may
therefore retain an arbitrary factor, or finitely many factors,
while bounding all other good background factors by
$2\delta$ on principal slots and by
$\delta e^{-wn}$ on animal slots.  Principal input weights are fixed
constant costs.  A fixed number of distinguished occurrences adds
only a fixed power of $k$, which a geometric series absorbs.  For
one retained factor the terms of order $k\ge2$ carry
$C_L^k\delta^{k-1}$; their sum is $O(C_L\delta)$ after enlarging
$C_L$ and taking $\delta$ smaller.  The same proof applies to finite
rectangular arrays of input indices.  Applying the multilinear
bounds to each elementary tensor in \eqref{block:projective}, summing
its products of norms, and taking the infimum proves the projective
statements.  For tensor outputs allocate the weights and counting
slack separately to the finitely many factors.  This explains why
only finitely many sufficiently separated weights were fixed.

Finally, the background connected series converges whenever
$C_L\delta<1$.  With finitely many square-free labels, at most that
many marked occurrences can appear; the background series still
converges and the resulting jet is polynomial in the bounded marked
coefficients.  Exponentiation of the absolutely convergent connected
series recovers the transformed moment series.  For finite lists this
identity first follows near zero from formal power series and extends
through the small-entry ball by convergence.  The same estimates
permit passage to absolutely summable slot lists.  Projection
\eqref{block:centering} costs at most two, and a saved scalar costs at
most the uncentered function norm, so all the bounds survive these
operations.
\end{proof}

\begin{corollary}[Large generating footprints]
\label{block:diameter}
Retain the animal footprint of a connected coefficient before projecting
onto constants.  For background coefficients from the small-entry
expansion whose footprint has physical diameter at least $D$, the total
absolute sum at a fixed output anchor is at most
\[
 C_L\delta e^{-cD/R}.
\]
The same gain can be retained together with any fixed number of
specified input norm factors, including marked factors.  On a torus of
side $N$, writing $m=N/R$, the total absolute sum of small-expansion
scalar terms whose footprints wrap or fail to fit in a chart of fixed
relative size is at most $C_L\delta m^2e^{-cm}$.  Here $m$ is sufficiently
large relative to the fixed halos.  Pure bad scalar terms have chart
footprints and do not contribute to this last sum.
\end{corollary}

\begin{proof}
A connected union of $n'$ parent cells has diameter at most $Cn'R$.
Run Lemma~\ref{block:tree} with twice the required output weight, and
retain one weight after the other has paid the ordinary norm.  The
remaining factor is at most $e^{-cD/R}$.  Every coefficient considered
contains a small entry, so its summed background bound is $C_L\delta$;
the sum over orders at least two is smaller.  A short singleton or pure
bad output has a footprint of diameter $CR$ and cannot be nonchart
when $m$ is large.  A nonchart connected footprint crosses a distance
$cN$ in torus coordinates, or has a lift crossing that distance in the
wrapping case, and therefore contains at least $cm$ parent cells.
There are at most $Cm^2$ anchors.  The preceding estimate, applied
before the scalar loses its footprint, gives the claimed total.
Retained norm factors survive because the tree proof bounds every
occurrence separately.
\end{proof}

\begin{corollary}[Uniform smallness of the tails]
\label{block:tails}
For sufficiently large $L$ and then sufficiently small $\delta$, the
transformed nonprincipal background satisfies
$\|h'_{A,c}\|_\infty\le\delta e^{-w|A|}$ for every disorder realization.
More precisely, with the same output weight $w$,
\begin{equation}
 \sup_{(A,c)}e^{w|A|}\|h'_{A,c}\|_\infty
       \le C_L\delta^2+C_Le^{-cL}\delta.
 \label{block:tail-propagation}
\end{equation}
\end{corollary}

\begin{proof}
Pure bad outputs are principal.  Short singleton outputs are also
principal, regardless of nearby bad fields.  Every remaining output
is therefore either of order at least two in small entries or comes
from a singleton animal of size $n\ge L/2$.  The first class is bounded
by $C_L\delta^2$ with weight $w$ by Lemma~\ref{block:tree}.
For the second, the weight allocation in
\eqref{block:compression} leaves a fixed positive fraction of $wn$
unused.  Since $n\ge L/2$, this gives the factor $e^{-cL}$ while
retaining enough slack for the placement sum.  Conditional averaging
under any bad tilt cannot enlarge its sup norm.  These two estimates
give \eqref{block:tail-propagation}.  Choose $L$ to make the second
coefficient at most $1/2$, and then $\delta$ to make
$C_L\delta\le1/2$.
\end{proof}

\subsection{Finite volumes and exact agreement with the plane}

\begin{proposition}[Finite-volume implementation]
\label{block:finite}
The preceding transformation and partition identity apply to finite
boxes with fixed or free boundary data, and to tori whose side is
an integer multiple of the working scales, while the scale is
sufficiently below the side length.  The input finite-volume lists
have finitely many possible anchors and admissible source jets.
On a torus, wrapping interactions remain in torus animal slots.  At
every fixed number of steps the principal norms have deterministic
bounds uniform in volume, disorder and decision seeds whenever the
microscopic principal norms have such a bound.  The total list sums
are finite, with bounds that may grow with the volume.  Bulk slots whose
full ancestral footprints avoid the boundary agree exactly with the
plane construction.  Infinite-plane lists require only the indexed
and rooted sums proved above; no infinite-volume partition integral
is used in their definition.
\end{proposition}

\begin{proof}
On a torus use torus distances and connected torus animals.  Make cuts
only in injective square charts of radius comparable with $HR$; the
scale is required to be smaller than a fixed sufficiently small
multiple of the torus side.  A wrapping carrier is retained as an
animal.  The counting and walk-covering arguments hold in charts,
and identifying possible placements can only reduce their number.
In a twisted torus the local reference specification includes the
twist, with a local gauge choice in each chart.  Only the background
deviation from that reference enters the list.  Center chart terms
after identifying that local reference with the plane reference.
A spin gauge change conjugates the conditional operators and
multiplication of functions, preserves sup norms and threshold tests,
and commutes with the geometric rules.  Thus, using the same decision
seeds, the complete local map is gauge covariant.  The saved scalars
of corresponding chart terms are identical because their reference
expectations agree after the gauge identification.  The possible
constant spin reversal between two gauges acts trivially on the
even background.  In particular,
these scalars cancel when comparing a torus and a twisted torus on
matching chart footprints.  The conditional identities and positivity
argument are unchanged.  The possible nonchart scalar discrepancy
is bounded by Corollary~\ref{block:diameter}.

For a box use an ambient grid extending past it and keep prescribed
exterior spins as formal coordinates until they are evaluated.  Forbid
a selected or ordinary singleton cut if its rule neighborhood meets
a fixed scale halo of the boundary.  Retain its insertions instead;
short entries keep the principal assignments already specified.
When testing singleton clearance, include every retained bad
insertion, including those retained because of the boundary rule.
All cuts that remain lie entirely in the reference bulk, so their
conditional kernels are exactly the reference kernels used above.
Center using the formal plane convention before evaluating exterior
spins, or use a fixed bounded projection off constants for boundary
entries.  No smallness of a boundary principal slot is needed.
The termwise proof of \eqref{block:tilted-identity} therefore applies
in each finite box.

There are only finitely many possible anchors at a fixed step in a
finite system: initially they lie in a bounded ambient grid, and each
step moves an anchor only to a parent or a bounded neighboring parent.
For each anchor the number of animals of size $n$ is at most
$e^{a_{\rm an}n}$.  Since $w>a_{\rm an}$, the tail bound proves total
absolute summability of all background animal slots.  The same
argument, with their fixed exponential weights and finitely many
mark anchors, proves absolute summability for source jets at every
fixed number of steps.  Their fixed factors $L^{qd}$ do not affect
this conclusion.  Truncating the animal lists and passing to the
limit proves the full finite-volume partition identity and justifies
all finite-volume expansions used in Proposition~\ref{block:exact}.

For completeness let $K$ bound the input principal norms at a given
step.  A selected component producing a fixed principal output has
at most $C_L$ fine centers; only $C_L$ such centers can contribute to
the same output halo.  Equation~\eqref{block:positive-log} bounds its
centered logarithm by $C_LK$.  Retained principal inputs obey the
same bound simply by summation.  The convergent small-entry series
adds at most a constant $C_L$.  Hence
\begin{equation}
 K'\le C_L(1+K).
 \label{block:finite-principal-bound}
\end{equation}
Induction gives a finite bound at each fixed step, independent of the
volume.  This bound need not be uniform in the number of steps and is
not a stochastic smallness assertion.

Finally, every output footprint includes the full input footprints,
all cut contours, all selection and threshold neighborhoods, and the
boundary-clearance tests.  If these footprints avoid the boundary,
the finite-volume and plane operations read exactly the same spin
functions, edge variables and seeds, so they give exactly the same
slot values.  Iterate this statement along the ancestors of a slot.
One may equivalently impose boundary clearance by a fixed multiple
of each current scale; the accumulated clearances are bounded by
$C\sum_{a\le j}r_a\le C'r_j$.  This proves exact bulk agreement and
completes the finite-volume construction.
\end{proof}

The transformation thus separates two tasks.  Positive conditional
averaging handles arbitrarily large principal backgrounds without a
loss depending exponentially on their norms.  The connected expansion
preserves small tails and records every source scalar and dependence
carrier.  We next estimate these exact lists in the disorder probability
space.
